
\section{Approximation-Based Orbit Acceleration}
\label{sec:approx-accel}
\crefname{lstlisting}{listing}{listings}
\crefformat{lstlisting}{#2listing~#1#3}

Perturbation theory (\cref{sec:perturbation-concept}) reduces the
precision required for per-pixel iteration, but it does not reduce the
\emph{number} of iterations: a pixel that does not escape still executes
the recurrence up to the requested limit~$n_{\max}$. At deep zooms that
limit can reach billions, and the work is repeated for every pixel.
The stored reference orbit has its own length~$n_{\mathrm{ref}}$;
rebasing permits evaluation to continue beyond that length. Neither
quantity is necessarily the orbit's period~$p$.

Approximation methods attack this remaining bottleneck by \emph{skipping}
large blocks of iterations entirely.  Instead of stepping through the
perturbation recurrence one iteration at a time, they precompute linear
maps that advance the perturbation delta by many iterations in a single
evaluation. \FractalShark{} implements two strategies:
\begin{itemize}
\item \textbf{Bilinear Approximation (BLA)}
  (\cref{sec:bilinear-approx})---precomputes composable two-term linear
  maps from the reference orbit, skipping iterations where the
  perturbations satisfy the step's validity test.
\item \textbf{Linear Approximation v2 (LA)}
  (\cref{sec:la-v2-perturb})---precomputes staged linear-approximation
  coefficients that can skip entire sections of the orbit in one
  evaluation, achieving larger skips at the cost of more complex
  coefficient management.
\end{itemize}
Both techniques preserve the perturbation framework's mixed-precision
advantage while replacing many of the $n_{\max}$ individual updates
with block evaluations. The achievable reduction depends on the orbit,
the pixel offset, and the validity thresholds.

\begin{figure}[!htbp]
\centering
\resizebox{\linewidth}{!}{%
\begin{tikzpicture}[
    >=Stealth,
    scale=0.75,
    dot/.style={circle, fill, inner sep=1.2pt},
    arr/.style={->, thick},
]
  % Common iteration axis ticks
  \foreach \x in {0,1,...,16} {
    \draw[gray!40] (\x, -0.1) -- (\x, 0.1);
  }
  \node[font=\scriptsize, below, gray] at (0, -0.15) {$0$};
  \node[font=\scriptsize, below, gray] at (4, -0.15) {$4$};
  \node[font=\scriptsize, below, gray] at (8, -0.15) {$8$};
  \node[font=\scriptsize, below, gray] at (12, -0.15) {$12$};
  \node[font=\scriptsize, below, gray] at (16, -0.15) {$16$};
  \draw[->, gray] (-0.3, 0) -- (17, 0)
    node[right, font=\scriptsize] {iteration $n$};

  % --- Row 1: Perturbation (step by step) ---
  \node[font=\scriptsize, anchor=east, orange!70!black] at (-0.5, 1.0)
    {Perturbation};
  \foreach \x in {0,1,...,15} {
    \pgfmathtruncatemacro{\xn}{\x+1}
    \draw[arr, orange!70!black, thin] (\x, 1.0) -- (\xn, 1.0);
  }
  \foreach \x in {0,1,...,16} {
    \node[dot, orange!70!black] at (\x, 1.0) {};
  }
  \node[orange!70!black, font=\scriptsize] at (8, 0.5)
    {16 individual updates};

  % --- Row 2: BLA (variable-length skips) ---
  \node[font=\scriptsize, anchor=east, green!50!black] at (-0.5, 2.8)
    {BLA};
  % Illustrative valid BLA spans.
  \draw[arr, green!50!black, thick, bend left=30] (0, 2.6) to (2, 2.6);
  \draw[arr, green!50!black, thick, bend left=30] (2, 2.6) to (4, 2.6);
  \draw[arr, green!50!black, thick, bend left=30] (4, 2.6) to (5, 2.6);
  \draw[arr, green!50!black, thick, bend left=30] (5, 2.6) to (8, 2.6);
  \draw[arr, green!50!black, thick, bend left=30] (8, 2.6) to (9, 2.6);
  \draw[arr, green!50!black, thick, bend left=30] (9, 2.6) to (12, 2.6);
  \draw[arr, green!50!black, thick, bend left=30] (12, 2.6) to (13, 2.6);
  \draw[arr, green!50!black, thick, bend left=30] (13, 2.6) to (16, 2.6);
  \foreach \x in {0,2,4,5,8,9,12,13,16} {
    \node[dot, green!50!black] at (\x, 2.6) {};
  }
  \node[green!50!black, font=\scriptsize] at (8, 2.0)
    {Illustrative variable-length skips};

  % --- Row 3: LA~v2 (adaptive staged skips) ---
  \node[font=\scriptsize, anchor=east, red!60!black] at (-0.5, 4.6)
    {LA~v2};
  % Stage 0 entries have variable spans.
  \draw[arr, red!60!black, thick, bend left=22] (0, 4.4) to (3, 4.4);
  \draw[arr, red!60!black, thick, bend left=22] (3, 4.4) to (8, 4.4);
  \draw[arr, red!60!black, thick, bend left=22] (8, 4.4) to (11, 4.4);
  \draw[arr, red!60!black, thick, bend left=22] (11, 4.4) to (16, 4.4);
  \foreach \x in {0,3,8,11,16} {
    \node[dot, red!60!black] at (\x, 4.4) {};
  }
  \node[red!60!black, font=\scriptsize] at (8, 3.8)
    {Stage 0: four valid block evaluations};

  % Higher stage annotation
  \draw[arr, red!60!black, very thick, bend left=15, dashed]
    (0, 5.9) to (8, 5.9);
  \draw[arr, red!60!black, very thick, bend left=15, dashed]
    (8, 5.9) to (16, 5.9);
  \foreach \x in {0,8,16} {
    \node[dot, red!60!black] at (\x, 5.9) {};
  }
  \node[red!60!black, font=\scriptsize] at (8, 5.45)
    {Stage 1: two composed block evaluations};
\end{tikzpicture}%
}% end resizebox
\caption[Iteration skipping: perturbation vs.\ BLA vs.\ LA]{%
  \sloppy
  Schematic iteration strategies over 16 original iterations.
  \emph{Perturbation} (orange) evaluates every update.
  \emph{BLA} (green) selects valid aligned linear maps.
  \emph{LA~v2} (red, solid) uses entries with adaptively chosen spans;
  higher stages (red, dashed) compose consecutive entries.
  The spans illustrate the hierarchy rather than a measured traversal
  or a guarantee that entries cover a particular period. A pixel uses
  a block only when its validity and remaining-budget checks pass.}
\label{fig:approx-skipping}
\end{figure}

\subsection{Bilinear approximation (BLA) for orbit deltas}
\label{sec:bilinear-approx}

Bilinear approximation~\cite{Imagina,Thompson2023BLA} precomputes composable linear
maps that advance the perturbation delta by multiple iterations at once,
avoiding per-step evaluation when the linearization error remains small
(\cref{fig:approx-skipping}).

\subsubsection{Reference orbit and delta formulation}
The perturbation framework (\cref{sec:perturbation-concept}) decomposes
the orbit of a nearby pixel \(c = c_\star + \Delta c\) into a
precomputed reference orbit \(z_n^\star\) and a delta
\(\Delta z_n = z_n - z_n^\star\).  The exact delta recurrence
(\cref{eq:perturb-exact}) is
\[
\Delta z_{n+1} = 2z_n^\star\,\Delta z_n + (\Delta z_n)^2 + \Delta c.
\]
When \(\Delta z_n\) remains small, the quadratic term can be neglected,
yielding a linearized update:
\begin{equation}
\Delta z_{n+1} \approx A_n\,\Delta z_n + B_n\,\Delta c,\qquad
A_n = 2z_n^\star,\quad B_n = 1.
\end{equation}
BLA generalizes this into precomputed multi-step maps that \emph{jump} multiple
iterations while maintaining an explicit validity bound.

\subsubsection{What \code{BLA<T>} stores}
A \code{BLA<T>} instance stores two complex coefficients:
\[
A = A_x + iA_y,\qquad B = B_x + iB_y,
\]
plus:
\begin{itemize}
  \item \code{r2}: a squared-radius validity bound used during lookup,
  \item \code{l}: the number of Mandelbrot iterations summarized by this step.
\end{itemize}
These objects exist to accelerate \emph{escape-time evaluation} by evolving
\(\Delta z\) cheaply for many pixels, then reconstructing \(z \approx z^\star +
\Delta z\) to perform bailout checks consistent with Mandelbrot rendering.

\subsubsection{Applying a step: complex multiply-add}
The method \code{getValue} (taking arguments
\code{Real\-Delta\-SubN}, \code{Imag\-Delta\-SubN},
\code{Real\-Delta\-Sub0}, \code{Imag\-Delta\-Sub0})
applies:
\[
\Delta z \leftarrow A\,\Delta z + B\,\Delta c,
\]
expanded into real arithmetic:
\begin{align}
\Re(\Delta z') &= A_x \Re(\Delta z) - A_y \Im(\Delta z) + B_x \Re(\Delta c) - B_y \Im(\Delta c), \\
\Im(\Delta z') &= A_x \Im(\Delta z) + A_y \Re(\Delta z) + B_x \Im(\Delta c) + B_y \Re(\Delta c).
\end{align}

\subsubsection{Composing steps to build longer jumps}
If a first step maps \(\Delta z \mapsto A_1\Delta z + B_1\Delta c\) and a second maps
\(\Delta z' \mapsto A_2\Delta z' + B_2\Delta c\), the composition is:
\[
A_{\text{new}} = A_2 A_1,\qquad B_{\text{new}} = A_2 B_1 + B_2.
\]
In the code (\code{BLA.h}), these two steps are passed as the parameters
\code{x} and \code{y} to \code{getNewA} and \code{getNewB}, corresponding to
step~1 and step~2 respectively.
Composition supports a hierarchy of step sizes (often powers of two) for quickly
advancing delta orbits while rendering the escape-time field.

\subsubsection{GPU lookup: selecting a valid aligned step}
A GPU-side helper such as \code{GPU\_BLAS} stores the hierarchy and selects a
step that is both \emph{aligned} with the current iteration index and
\emph{valid} under the current bound check (typically comparing a computed
squared-magnitude proxy \code{z2} against \code{r2}). When a step is valid, the
renderer can advance the orbit by \code{l} iterations at a cost far below
performing \code{l} individual perturbation updates.

\subsection{LA~v2 linear approximation with perturbation}
\label{sec:la-v2-perturb}

The LA~v2 pipeline~\cite{Imagina} combines staged approximation maps with
an exact perturbation recurrence for the remaining updates. Its purpose
is to amortize reference-dependent work across the image: the tables are
built once for a reference, while each pixel evaluates only the blocks
that remain valid for its state. Near a mini-Mandelbrot copy, returns
close to the critical point can provide useful block boundaries and
enable long skips.

CPU and GPU rendering both use
\code{FractalShark::LA::\allowbreak AdvancePixel} for AT initialization
and LA traversal. The CPU supplies an orbit reader backed by its
runtime decompressor; the GPU supplies one backed by
\code{GPUPerturbSingleResults}. The rendering pipeline has three
operations: optional Approximation Transform (AT), descent through
the LA hierarchy, and a perturbation tail. The AT and LA operations
advance an approximate state; the tail performs the rendering bailout
test. The perturbation-only case is described in \cref{sec:perturb-only}.

\subsubsection{Parameter delta per pixel}
Each pixel constructs a parameter offset \(\Delta c\) relative to a selected
reference center (sign conventions may incorporate the image \(Y\)-flip):
\begin{align}
\Delta c_x &= dx\cdot X - \code{centerX}, \\
\Delta c_y &= -dy\cdot Y - \code{centerY},
\end{align}
and packs \(\Delta c=\Delta c_x+i\Delta c_y\) into \code{DeltaSub0}.

\subsubsection{Delta state and reconstruction}
The renderer maintains a parameter offset \code{DeltaSub0}, a changing
orbit delta \code{DeltaSubN}, and an accumulated pixel iteration
count~$n$. Initially $n=0$ and $\Delta z=0$, unless AT advances the
pixel first. An absolute orbit estimate is reconstructed from the
current reference sample:
\[
z \approx z_r^\star + \Delta z.
\]
Here $r$ is a reference-orbit index, distinct from~$n$. Rebasing changes
$r$ without changing the accumulated count. During LA traversal, the
local index~$j$ selects an entry in a particular stage rather than an
individual orbit sample. A failed entry translates that index into the
next finer stage; at stage~0 the translation gives the raw orbit index
needed by the perturbation tail (\cref{sec:la-stage-hierarchy}).

\subsubsection{Perturbation update}
Given a reference sample \(z^\star=x^\star+iy^\star\) and \(\Delta z=\Delta
x+i\Delta y\), the perturbation form is:
\[
\Delta z \leftarrow \Delta z\cdot(2z^\star+\Delta z) + \Delta c.
\]
In real arithmetic, with temporaries corresponding to \((2x^\star+\Delta x)\)
and \((2y^\star+\Delta y)\), this yields:
\begin{align}
\Delta x' &= \Delta x\,(2x^\star+\Delta x) - \Delta y\,(2y^\star+\Delta y) + \Delta c_x, \\
\Delta y' &= \Delta x\,(2y^\star+\Delta y) + \Delta y\,(2x^\star+\Delta x) + \Delta c_y,
\end{align}
where both right-hand sides use the pre-update values of \(\Delta x\) and
\(\Delta y\) (the code saves these as \code{DeltaSubNXOrig} and
\code{DeltaSubNYOrig}).

\subsubsection{Escape test}
The perturbation tail reconstructs $z$ after each delta update and tests
its Euclidean squared magnitude. The rendering bailout in these LA~v2
paths is
\[
|z|^2=x^2+y^2,\qquad R_{\mathrm{bail}}^2=256.
\]
The CPU tail stops when $|z|^2>256$; the GPU tail continues only while
$|z|^2<256$, so equality also stops the GPU. Both tails test the newly
reconstructed value before increasing the stored iteration count for
that update. The mathematical Mandelbrot escape radius~2 corresponds
to squared magnitude~4; the renderer uses a later bailout at radius~16.
For HDR types, the norm is reduced before the specialized comparison.

LA traversal itself has no rendering bailout after a block. Its
reconstructed endpoint is used for rebasing, and a rejected block
causes descent to a finer stage. AT has a separate transformed-domain
limit (\cref{sec:attractor-transform}). These approximation checks
must be distinguished from the perturbation tail's escape test.

\subsubsection{LA coefficient model}
\label{sec:la-coeff-model}

The LA~v2 algorithm approximates the perturbation delta after multiple
iterations as a linear function of two quantities: a nonlinearly transformed
input perturbation and the parameter offset.  Concretely, each LA entry stores
a reference orbit value \(z_\mathrm{ref}\) together with two complex
coefficients \code{ZCoeff} and \code{CCoeff}. Its other fields describe
where the map is applicable and how traversal continues:
\begin{itemize}
  \item \code{LAThreshold} bounds the prepared perturbation;
        \code{LAThresholdC} bounds the parameter offset.
  \item \code{StepLength}, denoted $\ell$, counts the original
        Mandelbrot iterations covered by the entry.
  \item \code{NextStageLAIndex} identifies the same starting boundary
        in the next finer stage, or in the raw orbit for stage~0.
  \item \code{MinMag} is construction state for minimum-magnitude
        boundary detection; rendering does not use it.
\end{itemize}
Each stage stores its starting array offset \code{LAIndex} and its
number of evaluable entries \code{MacroItCount}. An additional terminal
entry stores the reference at the final endpoint. The reference of
entry~$j+1$, including that terminal entry, reconstructs the output of
entry~$j$; the terminal entry is not an extra block to evaluate.

Application proceeds in two stages:
\begin{enumerate}
  \item \textbf{Prepare} (nonlinear).  Given the current perturbation
        \(\Delta z\), compute
        \begin{equation}
        \tilde{\Delta z} = \Delta z\cdot(2\,z_\mathrm{ref} + \Delta z).
        \label{eq:la-prepare}
        \end{equation}
        Writing $s_x = 2\,z_{\mathrm{ref},x} + \Delta z_x$ and
        $s_y = 2\,z_{\mathrm{ref},y} + \Delta z_y$, the real-arithmetic
        expansion is
        \begin{align}
        \tilde{\Delta z}_x &= \Delta z_x\,s_x - \Delta z_y\,s_y, \\
        \tilde{\Delta z}_y &= \Delta z_x\,s_y + \Delta z_y\,s_x.
        \end{align}
        The operation retains the quadratic term at the entry's reference point and
        corresponds to one exact perturbation step (without the
        \(+\Delta c\) contribution).  The Prepare step also checks the
        validity condition
        \(\lVert\tilde{\Delta z}\rVert_\infty < \code{LAThreshold}\);
        if this fails the LA entry is unusable and the kernel falls back
        to a finer stage or to direct perturbation.

  \item \textbf{Evaluate} (linear).  The output perturbation after the
        iterations covered by this LA entry is
        \begin{equation}
        \Delta z_\mathrm{out}
          \approx \code{ZCoeff}\cdot\tilde{\Delta z}
          + \code{CCoeff}\cdot\Delta c.
        \label{eq:la-evaluate}
        \end{equation}
        In real arithmetic, with $Z_x,Z_y$ and $C_x,C_y$ denoting the components
        of \code{ZCoeff} and \code{CCoeff}:
        \begin{align}
        \Delta z_{\mathrm{out},x}
          &= Z_x\,\tilde{\Delta z}_x - Z_y\,\tilde{\Delta z}_y
           + C_x\,\Delta c_x - C_y\,\Delta c_y, \\
        \Delta z_{\mathrm{out},y}
          &= Z_x\,\tilde{\Delta z}_y + Z_y\,\tilde{\Delta z}_x
           + C_x\,\Delta c_y + C_y\,\Delta c_x.
        \end{align}
\end{enumerate}
The CPU \code{LAInfoDeep} and GPU \code{GPU\_LAInfoDeep} wrappers both
delegate Prepare and Evaluate to \code{FractalShark::LA} helpers.
For an entry spanning $[a,b)$, reconstruction uses
$z_b^\star+\Delta z_{\mathrm{out}}$, while Prepare uses $z_a^\star$.

The initial conditions, set in the \code{LAInfoDeep} constructor, are
\[
\begin{gathered}
\code{ZCoeff}_0 = 1,\quad \code{CCoeff}_0 = 1,
\\
\code{LAThreshold}_0 = 1,\quad \code{LAThresholdC}_0 = 1,
\\
\code{MinMag}_0 = 4\quad\text{for detection method 1}.
\end{gathered}
\]
The initial \(\code{CCoeff}_0 = 1\) accounts for the single \(+\Delta c\)
contribution implicit in combining the Prepare step with one perturbation
iteration: \(\Delta z_\mathrm{next} = \tilde{\Delta z} + \Delta c\), so
\(\Delta z_\mathrm{out} = 1\cdot\tilde{\Delta z} + 1\cdot\Delta c\).
The thresholds start at the finite bound~1 and tighten monotonically
as iterations are absorbed. The initial \code{MinMag} is a seed for a
running Chebyshev magnitude, not a squared norm or an escape test.
Method~2 uses threshold tightening instead and does not need that seed.

\subsubsection{Step-forward recurrence}
\label{sec:la-step-forward}

Each call to \code{LAInfoDeep::Step} absorbs one additional reference-orbit
iteration into an LA entry by appending a linearized perturbation step.
Starting from the exact perturbation recurrence (dropping the quadratic
\((\Delta z)^2\) term for subsequent iterations after Prepare):
\[
\Delta z_{n+1} \approx 2\,z_n^\star\,\Delta z_n + \Delta c,
\]
and substituting the LA model
\(\Delta z_n \approx \code{ZCoeff}_k\cdot\tilde{\Delta z}
                    + \code{CCoeff}_k\cdot\Delta c\), one obtains
\begin{align}
\Delta z_{n+1}
  &\approx 2\,z_n^\star\bigl(
      \code{ZCoeff}_k\cdot\tilde{\Delta z}
    + \code{CCoeff}_k\cdot\Delta c
    \bigr) + \Delta c \\
  &= \bigl(2\,z_n^\star\cdot\code{ZCoeff}_k\bigr)\,\tilde{\Delta z}
   + \bigl(2\,z_n^\star\cdot\code{CCoeff}_k + 1\bigr)\,\Delta c.
\end{align}
Reading off the coefficients gives the step-forward recurrence:
\begin{equation}
\code{ZCoeff}_{k+1} = 2\,z_n^\star\cdot\code{ZCoeff}_k,\qquad
\code{CCoeff}_{k+1} = 2\,z_n^\star\cdot\code{CCoeff}_k + 1.
\label{eq:la-step}
\end{equation}

Expanding each complex product $2\,z_n^\star \cdot \code{ZCoeff}_k$
(and similarly for \code{CCoeff}) into components:
\begin{align}
Z_{x,k+1} &= 2\,(x_n^\star\,Z_{x,k} - y_n^\star\,Z_{y,k}), \\
Z_{y,k+1} &= 2\,(x_n^\star\,Z_{y,k} + y_n^\star\,Z_{x,k}), \\
C_{x,k+1} &= 2\,(x_n^\star\,C_{x,k} - y_n^\star\,C_{y,k}) + 1, \\
C_{y,k+1} &= 2\,(x_n^\star\,C_{y,k} + y_n^\star\,C_{x,k}).
\end{align}
The $+1$ contributes only to the real part of \code{CCoeff}, reflecting
the direct $+\Delta c$ term in the Mandelbrot recurrence.

\paragraph{Connection to the orbit derivative.}
The \code{CCoeff} recurrence
\(C_{k+1} = 2\,z_n^\star\,C_k + 1\) with \(C_0 = 1\)
has the same form as the orbit derivative recurrence
\(d_{n+1} = 2\,z_n\,d_n + 1\) with \(d_0 = 0\)
(\cref{subsec:ref-orbit-periodicity}), offset by one step.
The coefficient accumulates the parameter contributions introduced
\emph{within the block}. For a block starting at the critical point,
it agrees with the reference orbit's parameter derivative at the block
endpoint. For a general input whose state already depends on~$c$,
the incoming derivative must also propagate through Prepare and
\code{ZCoeff}; \code{CCoeff} alone is not the total derivative
(\cref{sec:la-derivatives}).

\paragraph{Worked three-iteration map.}
For $c_\star=-1/2$, the reference samples are
$z_0^\star=0$, $z_1^\star=-1/2$, $z_2^\star=-1/4$, and
$z_3^\star=-7/16$. A block starting at $z_0^\star$ has the following
coefficient development:
\begin{center}
\begin{tabular}{clcc}
\hline
Span & Operation & \code{ZCoeff} & \code{CCoeff} \\
\hline
1 & Prepare and seed & $1$ & $1$ \\
2 & Append $z_1^\star=-1/2$ & $-1$ & $0$ \\
3 & Append $z_2^\star=-1/4$ & $1/2$ & $1$ \\
\hline
\end{tabular}
\end{center}
Writing the incoming delta as $\xi$ and $\Delta c=\varepsilon$,
Prepare gives $\tilde{\Delta z}=\xi^2$, and the block estimates
\[
z_3\approx -\frac{7}{16}+\frac{1}{2}\xi^2+\varepsilon.
\]
The incoming quadratic term survives even though the remaining
updates are linearized. For a fresh critical orbit, $\xi=0$, exact
iteration instead gives
\[
z_3=-\frac{7}{16}+\varepsilon-\frac{1}{2}\varepsilon^2
    +\varepsilon^4.
\]
The omitted terms show why a validity test is needed even when the
first update is exact. The example illustrates coefficient algebra;
it does not prescribe the boundaries chosen by table construction.

\subsubsection{Validity threshold}
\label{sec:la-threshold}

The linearization underlying the step-forward recurrence introduces error
whenever the neglected quadratic term \((\Delta z)^2\) is non-negligible.
To control this, each LA entry tracks a validity threshold
\code{LAThreshold} that bounds how large the prepared perturbation
\(\tilde{\Delta z}\) may be.  Each step tightens the bound:
\begin{equation}
\code{LAThreshold}_{k+1}
  = \min\!\bigl(
      \code{LAThreshold}_k,\;
      \lVert z^\star\rVert_\infty \,/\,
      \lVert\code{ZCoeff}_k\rVert_\infty
      \;\cdot\;\code{LAThresholdScale}
    \bigr),
\label{eq:la-threshold}
\end{equation}
where \(\lVert\cdot\rVert_\infty\) is the Chebyshev (max-component) norm,
\(\lVert a + bi \rVert_\infty = \max(|a|, |b|)\), and
\code{LAThresholdScale} is a tunable safety factor, with default
$2^{-24}$ in \code{LAParameters}. At an appended sample, the neglected
term has magnitude $|\Delta z|^2$, whereas the linear term has
magnitude $|2z^\star\Delta z|$. Keeping the amplified perturbation
small relative to the local reference magnitude limits the relative
importance of the quadratic term. Division by the current coefficient
magnitude translates that local restriction back to the block input.
The minimum retains every restriction already encountered.

These tests use max-component magnitudes and configurable margins;
they are practical controls on the approximation domain, not a
certified bound on the final pixel error. Larger margins admit more
work to the approximation path but reduce the separation between the
retained linear term and the omitted nonlinear terms.

\paragraph{Parameter-offset threshold (\code{LAThresholdC}).}
A companion threshold \code{LAThresholdC} controls validity with respect
to the parameter offset \(\Delta c\).  It follows the same monotonically
decreasing pattern but is normalized by \(\lVert\code{CCoeff}\rVert\):
\begin{equation}
\code{LAThresholdC}_{k+1}
  = \min\!\bigl(
      \code{LAThresholdC}_k,\;
      \lVert z^\star\rVert_\infty \,/\,
      \lVert\code{CCoeff}_k\rVert_\infty
      \;\cdot\;\code{LAThresholdCScale}
    \bigr).
\label{eq:la-threshold-c}
\end{equation}
Its default scale is also $2^{-24}$. Before entering a stage's LA loop,
the shared traversal checks whether the pixel's
parameter offset satisfies the validity bound \eqref{eq:la-threshold-c}:
\[
\lVert \Delta c \rVert_\infty \ge \code{LAThresholdC}
  \quad\Longrightarrow\quad
  \text{skip this stage entirely.}
\]
The gate uses \code{LAThresholdC} from the \emph{first entry} of the
stage. Parameter validity and prepared-delta validity are separate
conditions: a small prepared delta does not compensate for an excessive
parameter offset. Equality rejects the stage or prepared delta.
Threshold scales are constructed from their configured exponents in
the generation scalar type by \code{LAParametersRuntime}, preserving
that type's exponent range.

\subsubsection{Composing LA entries}
\label{sec:la-composite}

Two consecutive LA entries can be merged into a single entry covering their
combined iteration span.  If LA\(_1\) covers iterations \([a,b)\) with
coefficients \((\code{ZCoeff}_1, \code{CCoeff}_1)\) and reference
\(z_\mathrm{ref,1}\), and LA\(_2\) covers \([b,c)\) with
\((\code{ZCoeff}_2, \code{CCoeff}_2)\) and reference \(z_\mathrm{ref,2}\),
then composing them requires one linearized Prepare step at the boundary
(dropping the quadratic term):
\begin{align}
\code{ZCoeff}_\mathrm{comp}
  &= 2\,z_\mathrm{ref,2}\cdot\code{ZCoeff}_1\cdot\code{ZCoeff}_2,
  \label{eq:la-comp-z} \\
\code{CCoeff}_\mathrm{comp}
  &= 2\,z_\mathrm{ref,2}\cdot\code{CCoeff}_1\cdot\code{ZCoeff}_2
   + \code{CCoeff}_2.
  \label{eq:la-comp-c}
\end{align}
In real arithmetic, the composition proceeds via intermediate products.
Let $Z' = 2\,z_{\mathrm{ref},2} \cdot \code{ZCoeff}_1$ (and similarly
$C' = 2\,z_{\mathrm{ref},2} \cdot \code{CCoeff}_1$); then:
\begin{align}
Z'_x &= 2\,(z_{\mathrm{ref},2,x}\,Z_{1,x} - z_{\mathrm{ref},2,y}\,Z_{1,y}), &
Z'_y &= 2\,(z_{\mathrm{ref},2,x}\,Z_{1,y} + z_{\mathrm{ref},2,y}\,Z_{1,x}), \\
Z_{\mathrm{comp},x} &= Z'_x\,Z_{2,x} - Z'_y\,Z_{2,y}, &
Z_{\mathrm{comp},y} &= Z'_x\,Z_{2,y} + Z'_y\,Z_{2,x},
\end{align}
and the composite \code{CCoeff} follows the same pattern with
$\code{CCoeff}_2$ added after the chained multiply.

The structure is analogous to BLA composition (\cref{sec:bilinear-approx}):
the \code{ZCoeff} factors multiply and the \code{CCoeff} factors compose
as an affine chain.  The factor \(2\,z_\mathrm{ref,2}\) accounts for
the linearized Prepare step at the junction between the two entries.
The composite's reference is set to \(z_\mathrm{ref,1}\), since the
composite's Prepare step operates at the first entry's starting point.

\paragraph{Threshold propagation during composition.}
The composite thresholds are \emph{not} the minimum of the two
entries' thresholds.  The computation proceeds in two stages.  First,
one Step-like tightening is applied at the junction point:
\begin{align}
T_\mathrm{junc}
  &= \min\!\bigl(
       \code{LAThreshold}_1,\;
       \lVert z_\mathrm{ref,2} \rVert_\infty
       \,/\, \lVert \code{ZCoeff}_1 \rVert_\infty
       \cdot \code{LAThresholdScale}
     \bigr),
\label{eq:la-comp-thresh-junc}
\\
T_{C,\mathrm{junc}}
  &= \min\!\bigl(
       \code{LAThresholdC}_1,\;
       \lVert z_\mathrm{ref,2} \rVert_\infty
       \,/\, \lVert \code{CCoeff}_1 \rVert_\infty
       \cdot \code{LAThresholdCScale}
     \bigr).
\end{align}

Then the intermediate product
\(Z' = 2\,z_\mathrm{ref,2} \cdot \code{ZCoeff}_1\) is formed
(this product is the ZCoeff \emph{before} multiplication by \code{ZCoeff}\(_2\)).
The second entry's threshold is rescaled by this amplification:
\begin{align}
\code{LAThreshold}_\mathrm{comp}
  &= \min\!\bigl(
       T_\mathrm{junc},\;
       \code{LAThreshold}_2 \,/\, \lVert Z' \rVert_\infty
     \bigr),
\\
\code{LAThresholdC}_\mathrm{comp}
  &= \min\!\bigl(
       T_{C,\mathrm{junc}},\;
       \code{LAThreshold}_2 \,/\, \lVert C' \rVert_\infty
     \bigr),
\end{align}
where \(C' = 2\,z_\mathrm{ref,2} \cdot \code{CCoeff}_1\).
The rescaling by \(\lVert Z' \rVert_\infty\) accounts for the fact that
the prepared perturbation entering LA\(_2\) is amplified by the first
entry's coefficients; the second entry's threshold must be correspondingly
tightened.
Both propagated restrictions use the second entry's
\code{LAThreshold}, since they limit the two contributions to that
entry's prepared input.

In the code, \code{LAInfoDeep::Composite} implements this operation and
is used by \code{CreateNewLAStage} to build higher-level stages from
lower-level ones.

\subsubsection{Adaptive stage hierarchy}
\label{sec:la-stage-hierarchy}

The hierarchy uses returns toward the critical point to choose useful
block boundaries, but its entries have adaptive lengths. The
construction variable named \code{Period} can denote a detected return
span or a synthetic target span; it need not be a proven dynamical
period. Near a periodic mini-copy, such boundaries often follow the
orbit's return structure. Validity thresholds still determine whether
a particular pixel can use those blocks.

\paragraph{Boundary detection.}
Stage~0 appends raw orbit updates with \code{LAInfoDeep::Step}.
Method~1, the default, tracks a minimum magnitude $m$, initially~4,
over the appended reference samples. For a candidate sample $z_k^\star$,
the update and boundary test are
\begin{equation}
m'=\min(m,\lVert z_k^\star\rVert_\infty),\qquad
m'<m\,\tau_0,
\label{eq:period-detect-minmag}
\end{equation}
where $\tau_0=\code{Stage0PeriodDetectionThreshold2}$, with default
$2^{-6}$. The starting reference value is stored separately; the
critical value $z_0^\star=0$ is not included in this running minimum.
A substantial drop identifies a useful return toward the critical
point rather than certifying an exact period.

Method~2 tests a sharp tightening of the prepared-delta threshold:
\begin{equation}
T'=\min\!\left(T,
  \frac{\lVert z_k^\star\rVert_\infty}
       {\lVert\code{ZCoeff}\rVert_\infty}
  \code{LAThresholdScale}\right),\qquad
T'<T\,\tau_0,
\label{eq:period-detect-coeff}
\end{equation}
where $T=\code{LAThreshold}$ and the method's $\tau_0$ is
\code{Stage0PeriodDetectionThreshold}, with default $2^{-10}$.
At higher stages the Composite operation detects boundaries at its
junction: method~1 tests the minimum including the junction reference,
and method~2 tests the junction tightening from
\cref{eq:la-comp-thresh-junc}, before the second entry's propagated
restriction. The corresponding higher-stage factors are
\code{PeriodDetectionThreshold2} (default $2^{-3}$) and
\code{PeriodDetectionThreshold} (default $2^{-10}$).
\code{DetectPeriod} also tests a prospective sample when deciding
whether to seed the next block with one or two updates.

\paragraph{Building stage~0.}
Construction seeds a block at $z_0^\star=0$ and appends $z_1^\star$,
giving an initial two-iteration map. If that map's \code{ZCoeff} is
zero, initialization fails. Otherwise, the scan continues until a
boundary is detected or the reference ends. When a candidate triggers
detection at index~$i$, the stored entry is the accumulated map
\emph{before} that candidate is absorbed. Its endpoint is $z_i^\star$;
the candidate belongs to the next block.

If the scan detects no boundary and $n_{\mathrm{ref}}\le64$, it stores
a single block and terminal reference but leaves the LA reference
invalid for rendering. Longer references receive a synthetic target
span. A detected initial span greater than~64 also causes a restart
with synthetic segmentation. Once the target is established, each
block ends either at a detected boundary or when that target span has
been reached. A final partial block covers the remainder.

\paragraph{Synthetic spans and higher stages.}
\code{CalculatePeriod} computes a target span using
\begin{equation}
k=\operatorname{round}\!\left(
  \frac{\log_2 n_{\mathrm{ref}}}{d}\right),\qquad
h=s\,\operatorname{round}\!\left(\rho^{1/k}\right),
\label{eq:la-synthetic-span}
\end{equation}
where $d=2$ for uncompressed references and $d=8$ for simple
compression. The larger compression divisor generally gives larger
target spans and fewer entries, trading finer fallback granularity
for reduced table memory. At stage~0, $s=1$ and $\rho=n_{\mathrm{ref}}$ for both
the no-boundary fallback and the oversized-span restart. At a higher
stage, $s$ is the first entry's step length in the preceding stage.
The fallback uses $\rho=n_{\mathrm{ref}}/s$; subdivision of an
oversized detected span~$h_{\mathrm{det}}$ uses
$\rho=h_{\mathrm{det}}/s$. Both roundings in
\eqref{eq:la-synthetic-span} are part of the construction rule.

\code{CreateNewLAStage} composes consecutive entries from the
preceding stage. The step lengths add, while the new
\code{NextStageLAIndex} records the first constituent's index in that
preceding stage. Boundaries again follow detection or the target span.
Without a detected boundary, synthetic grouping continues when
$n_{\mathrm{ref}}>64s$; otherwise the final stage contains a single
block and construction stops. Detected spans greater than $64s$ are
subdivided. The limits control grouping scale rather than imposing a
maximum number of entries on every stage.

The construction outline in \cref{lst:la-construction} uses descriptive operations for the
scan and segmentation; Step and Composite implement the coefficient
and threshold recurrences already derived.

\begin{lstlisting}[float=htbp,caption={LA table construction outline},label={lst:la-construction}]
mark reference invalid; disable AT
if reference has no iterations: return
seed map at orbit[0]; append orbit[1]
if ZCoeff is zero: return
scan raw samples to find initial boundary
if no boundary and reference length <= 64:
    store single block and terminal reference; return
if no boundary or initial span > 64:
    restart with synthetic stage-0 target span
segment raw orbit using Step:
    end block before a detected boundary or at target span
    store its length and raw starting index
    reseed next block; retain final partial block
append terminal reference; record stage offset and count
repeat:
    compose preceding-stage entries into the next stage
    detect boundaries or choose synthetic grouping
    store summed lengths and finer-stage starting indices
    append terminal reference; record stage offset and count
    stop when construction requests no further stage
select first usable AT from coarsest stage downward
mark reference valid
\end{lstlisting}

Stage~0 generation can use one thread or multiple threads. The
multithreaded path shares the initial span selection, finds compatible
starting boundaries for worker ranges, and assembles contiguous,
non-overlapping entry sequences before adding the terminal reference.
Small workloads use the single-threaded path. Higher stages are built
from the resulting stage~0 array by the same composition operation.

\paragraph{Descent, fallback, and rebasing.}
Evaluation starts at the coarsest stage. A first-entry parameter gate
can skip that stage entirely. Otherwise \code{getLA} rejects a block
if its full step length exceeds the remaining iteration budget or
its prepared delta reaches the validity threshold. Neither failure
partially applies the block. Instead, its \code{NextStageLAIndex}
becomes the entry index for the next finer stage. At stage~0 the same
field gives the raw reference index for the perturbation tail.

After a successful block, the renderer increases $n$ by $\ell$,
reconstructs $z$ using the next entry's reference, and advances~$j$.
If $\lVert\Delta z\rVert_\infty>\lVert z\rVert_\infty$ or $j$ reaches
\code{MacroItCount}, it sets $\Delta z=z$ and $j=0$. Every stage begins
at the zero reference, so the same absolute state can be represented
relative to that origin. The accumulated iteration count is preserved.
The perturbation tail uses the analogous magnitude and reference-end
rebasing rules, with squared Euclidean norms.

\begin{lstlisting}[float=htbp,caption={Shared per-pixel AT and LA traversal},label={lst:la-traversal}]
(n, deltaZ) = InitializePixel(deltaC, iterationLimit)
referenceIndex = 0
z = deltaC  // initial returned value if no step was taken
if n != 0 and an origin reference is available:
    z = orbit[referenceIndex] + deltaZ
for stage from coarsest to finest, if reference is valid:
    if normInfinity(deltaC) >= stage.first.LAThresholdC:
        continue
    j = referenceIndex
    while n < iterationLimit:
        entry = stage[j]
        if n + entry.StepLength > iterationLimit:
            referenceIndex = entry.NextStageLAIndex
            break
        prepared = deltaZ * (2 * entry.Ref + deltaZ)
        if normInfinity(prepared) >= entry.LAThreshold:
            referenceIndex = entry.NextStageLAIndex
            break
        n += entry.StepLength
        deltaZ = prepared * entry.ZCoeff
                 + deltaC * entry.CCoeff
        z = stage[j + 1].Ref + deltaZ
        ++j
        if normInfinity(deltaZ) > normInfinity(z)
           or j >= stage.MacroItCount:
            deltaZ = z
            j = 0
    if n >= iterationLimit: break
return (n, referenceIndex, deltaZ, z)
\end{lstlisting}

The traversal in \cref{lst:la-traversal} suppresses arithmetic reductions and storage offsets.
\code{AdvancePixel} initializes its returned $z$ to $\Delta c$ when
no approximation step has completed; the perturbation tail then
reconstructs the first updated absolute value. After AT advances the
pixel, the reader reconstructs at reference index~0; a supplied period
also permits origin reconstruction in the zero-length-reference
fallback. Stage gates do not advance the orbit state. At a rejected
entry, the finer-stage mapping preserves all completed work.

\paragraph{Worked index and budget trace.}
Consider an illustrative stage~0 with spans $2,3,2,1$, covering raw
boundaries $0,2,5,7,8$. A coarser stage composes the first two and the
last two entries into spans $5,3$. Its finer-stage indices are $0,2$,
respectively (\cref{fig:la-stage-fallback}). These metadata illustrate
indexing independently of any particular reference's detected spans.

\begin{figure}[!htbp]
\centering
\begin{tikzpicture}[>=Stealth, x=1.1cm, y=1cm,
    every node/.style={font=\small}]
  \draw[->, gray] (0,0) -- (8.5,0);
  \foreach \x in {0,2,5,7,8} {
    \draw (\x,-0.08) -- (\x,0.08);
    \node[below] at (\x,-0.08) {$\x$};
  }
  \node[anchor=east] at (-0.2,0) {Raw index};
  \node[anchor=east] at (-0.2,0.9) {Stage 0};
  \node[anchor=east] at (-0.2,2.1) {Stage 1};
  \foreach \a/\b/\j in {0/2/0,2/5/1,5/7/2,7/8/3} {
    \draw[fill=blue!8] (\a,0.55) rectangle (\b,1.25);
    \node at ({(\a+\b)/2},0.9) {$j=\j$};
  }
  \foreach \a/\b/\j in {0/5/0,5/8/1} {
    \draw[fill=red!8] (\a,1.75) rectangle (\b,2.45);
    \node at ({(\a+\b)/2},2.1) {$j=\j$};
  }
  \draw[->, thick] (0.3,1.75) -- (0.3,1.25);
  \draw[->, thick] (5.3,1.75) -- (5.3,1.25);
  \node[above] at (2.5,2.45) {$\ell=5$, finer index $0$};
  \node[above] at (6.5,2.45) {$\ell=3$, finer index $2$};
  \draw[dashed, orange!80!black] (6,-0.25) -- (6,2.85);
  \node[below, orange!80!black] at (6,-0.55) {$n_{\max}=6$};
\end{tikzpicture}
\caption[LA finer-stage index mapping]{Illustrative LA entry spans and
  finer-stage mappings. After the first coarse block advances to~5,
  the three-iteration coarse block and two-iteration fine block both
  exceed a six-iteration budget. Fallback reaches raw reference
  index~5, where one perturbation update can finish the budget.}
\label{fig:la-stage-fallback}
\end{figure}

Assume AT is disabled, both stage gates pass, all Prepare checks pass,
and no magnitude rebase occurs. With $n_{\max}=6$, the first coarse
block advances $n$ from~0 to~5. The next coarse block would advance
to~8, so it is rejected and traversal descends to fine entry~2.
That entry would advance to~7 and is also rejected. Its raw starting
index is~5, so Full rendering continues from that sample with one
perturbation update, provided the bailout test passes. With a budget
of~7, the fine block fits exactly and traversal finishes at~7.

\paragraph{Rendering modes.}
The CPU LA~v2 renderer uses the full pipeline. The CUDA entry point
\code{mandel\_1xHDR\_float\_perturb\_lav2} selects phases at compile time:
\begin{center}
\begin{tabular}{lll}
\hline
\code{LAv2Mode} & AT and LA traversal & Perturbation tail \\
\hline
\code{Full} & Enabled when applicable & Enabled \\
\code{LAO} & Enabled when applicable & Omitted \\
\code{PO} & Omitted & Enabled \\
\hline
\end{tabular}
\end{center}
LAO writes the count reached by approximation alone. When approximation
stops short, that count does not represent completion of the requested
escape-time evaluation. PO initializes a zero delta and performs every
update through the perturbation recurrence.

\subsubsection{Approximation Transform (AT)}
\label{sec:attractor-transform}

When a stage's first block defines a usable quadratic return map, the
\emph{Approximation Transform} evaluates that map repeatedly without
reference-orbit lookups. Such a map can be useful near a mini-Mandelbrot
copy, where many original iterations describe one local return. The
LA-derived approximation makes that local return inexpensive to
evaluate, subject to its transformed-domain and parameter checks.

\paragraph{Deriving the transformed recurrence.}
The first entry starts at $z_0^\star=0$, so Prepare is simply $z^2$.
Let its coefficients be $A=\code{ZCoeff}$ and $B=\code{CCoeff}$, its
step length be $\ell$, and its endpoint reference be
$q=z_{\mathrm{ref,next}}$. Combining Evaluate with reconstruction gives
\begin{equation}
z_{\mathrm{out}}\approx A z^2+B\Delta c+q.
\label{eq:at-return-map}
\end{equation}
With $w=Az$, multiplying the map by $A$ gives
\begin{equation}
w_{i+1}\approx w_i^2+c_w,\qquad c_w=AB\Delta c+Aq.
\label{eq:at-recurrence}
\end{equation}
The implemented recurrence evaluates the right-hand side directly.
Each transformed update accounts for $\ell$ original iterations;
$i$ counts transformed updates, while $n=i\ell$ counts original ones.
The selected $\ell$ can reflect a return span or a composed synthetic
span and need not equal a fundamental period.

\paragraph{Construction.}
The AT is constructed from the first entry of a stage's LA array.
The algorithm iterates from the topmost (coarsest) stage downward,
attempting to construct an AT at each level.  The first stage whose
AT has positive step length and passes the usability check
(\cref{sec:at-usability}) is
selected; if no stage produces a usable AT, the optimization is
disabled for this reference.  For the selected stage, let the first
entry have coefficients \code{ZCoeff}, \code{CCoeff} and let the
\emph{next} entry's reference be \(z_\mathrm{ref,next}\).  The
transformation parameters are:
\begin{align}
\code{AT.ZCoeff}   &= \code{ZCoeff}, \\
\code{AT.CCoeff}   &= \code{ZCoeff} \cdot \code{CCoeff}, \\
\code{AT.InvZCoeff}&= 1 \,/\, \code{AT.ZCoeff}, \\
\code{AT.RefC}     &= z_\mathrm{ref,next} \cdot \code{ZCoeff}, \\
\code{AT.SqrEscapeRadius} &= \min\bigl(
      \lVert\code{ZCoeff}\rVert^2 \cdot \code{LAThreshold},\;
      \Lambda\bigr), \\
\code{AT.ThresholdC} &= \min\bigl(
      \code{LAThresholdC},\;
      \Lambda \,/\, \lVert\code{AT.CCoeff}\rVert_\infty\bigr).
\end{align}
Here $\lVert\cdot\rVert^2$ is the Euclidean squared norm and
$\lVert\cdot\rVert_\infty$ is the Chebyshev norm. The cap is
\[
\Lambda=\begin{cases}
2^{256}, & \text{HDR double with small exponents disabled},\\
2^{32}, & \text{otherwise}.
\end{cases}
\]
\code{UseSmallExponents} is enabled when generating tables for a
2x32 target. The coefficients and transformed recurrence use the
selected scalar family, including HDR arithmetic when that family
requires it (\cref{sec:la-dp-conversion}).

The transformed radius has a validity interpretation. Since
$|w|^2=|A|^2|z|^2$ and
$\lVert z^2\rVert_\infty\le|z|^2$, keeping $|w|^2$ below
$|A|^2\code{LAThreshold}$ keeps the prepared value within the first
entry's domain. The cap limits the permitted transformed magnitude.
The parameter cap similarly restricts the amplification of $\Delta c$.

\paragraph{Usability check.}
\label{sec:at-usability}
For a candidate with positive step length, \code{ATInfo::Usable}
requires two strict inequalities:
\begin{enumerate}
  \item The reference's maximum perturbation radius $R_{\max}$ gives
        sufficient parameter variation relative to the constant term:
        \[
        R_{\max}^2 \cdot \lVert\code{AT.CCoeff}\rVert^2 \cdot 2^{32}
          > \lVert\code{AT.RefC}\rVert^2.
        \]
  \item The squared transformed radius must exceed~4:
        \[
        \code{AT.SqrEscapeRadius}>4.
        \]
\end{enumerate}
The first inequality is a reference-level selection heuristic using
the stored \code{factor}, whose default is $2^{32}$. It is not the
per-pixel validity test. Even for a selected AT, each pixel must satisfy
\[
\lVert\Delta c\rVert_\infty\le\code{AT.ThresholdC}.
\]
Unlike the LA stage and Prepare gates, this test accepts equality.
If the reference is invalid, AT is disabled, the step length is zero,
or the pixel fails this test, initialization returns $n=0$ and the
caller's zero delta for ordinary LA traversal.

\paragraph{Evaluation.}
\code{ATInfo::InitializePixel} owns the initialization operation for
both CPU and GPU rendering. For an eligible pixel it computes
\[
c_w = \Delta c \cdot \code{AT.CCoeff} + \code{AT.RefC},
\]
and starts at $w_0=0$. The transformed budget is
$i_{\max}=\lfloor n_{\max}/\ell\rfloor$. Before each update, the
operation tests whether $|w_i|^2>\code{AT.SqrEscapeRadius}$; equality
permits another update. No reference-orbit lookups occur in this loop.

In real arithmetic, the forward and inverse AT transforms expand as standard
complex multiply-adds.  Writing \code{AT.CCoeff} $= P_x + i\,P_y$ and
\code{AT.RefC} $= R_x + i\,R_y$:
\begin{align}
c_{w,x} &= \Delta c_x\,P_x - \Delta c_y\,P_y + R_x, \\
c_{w,y} &= \Delta c_x\,P_y + \Delta c_y\,P_x + R_y.
\end{align}
The transformed iteration expands as
$w'_x=w_x^2-w_y^2+c_{w,x}$ and
$w'_y=2w_xw_y+c_{w,y}$, with both right-hand sides using the old
components. The inverse transform
$\Delta z = w \cdot \code{AT.InvZCoeff}$ is another complex product.

After the domain test stops the loop or the transformed budget is
exhausted, initialization returns
\[
\Delta z=w_i\cdot\code{AT.InvZCoeff},\qquad n=i\ell,
\]
relative to the zero reference. The pseudocode in
\cref{lst:at-initialization} makes the test order and returned count explicit.

\begin{lstlisting}[float=htbp,caption={AT pixel initialization},label={lst:at-initialization}]
n = 0
deltaZ = callerZero
if not enabled or StepLength == 0:
    return (n, deltaZ)
if normInfinity(deltaC) > ThresholdC:
    return (n, deltaZ)
limit = iterationLimit / StepLength  // integer division
cw = deltaC * CCoeff + RefC
w = 0
i = 0
while i < limit:
    if normSquared(w) > SqrEscapeRadius: break
    w = w * w + cw
    ++i
deltaZ = w * InvZCoeff
n = i * StepLength
return (n, deltaZ)
\end{lstlisting}

If $n_{\max}<\ell$, no transformed update fits. If the budget is not
a multiple of $\ell$, the returned count leaves its remainder for
finer LA entries or the perturbation tail. Leaving the transformed
domain is a handoff, not a final rendering escape decision: the
original orbit magnitude and iteration count are resolved by the
remaining pipeline. A fully exhausted original budget requires no
additional LA work.

As an arithmetic trace, take $c_w=1$, $\ell=5$, an identity inverse
transform, an original iteration budget of~100, and transformed
squared-radius limit~256. The transformed
values are $0,1,2,5,26$. The next pre-update test observes
$26^2>256$, so initialization returns $i=4$, $n=20$, and
$\Delta z=26$. The values illustrate evaluation of supplied AT
parameters; candidate construction and reference usability remain
separate requirements.

\paragraph{Significance.}
When the map remains applicable, AT replaces roughly $\ell$ original
updates with one complex square and addition, while eliminating
reference traffic during that phase. The benefit grows with the
selected block length and the fraction of the pixel's work performed
before handoff. The LA-derived map explains the acceleration; the
usability and per-pixel tests determine where it is attempted.

\subsubsection{Arithmetic types and table conversion}
\label{sec:la-dp-conversion}

The selected rendering algorithm fixes the arithmetic family for its
LA entries, AT parameters, delta state, and perturbation tail. GPU
variants use float, double, 2x32, or their HDR forms; the CPU LA~v2
paths use HDR float or HDR double. HDR extends the exponent range
while retaining the selected significand precision
(\cref{sec:hdr-float}).

Precomputed data can be converted when preparing a reference for a
different target type. In particular, 2x32 rendering can generate
reference and LA data in the corresponding double family and convert
them for evaluation. Conversion preserves entry spans, stage offsets,
and finer-stage indices while converting numerical fields to the
destination type. Every stage evaluates in the destination arithmetic
family, and AT uses that family with the construction cap appropriate
to its target.

\subsubsection{Derivative propagation through LA steps}
\label{sec:la-derivatives}

LA coefficients also support derivative propagation. Feature Finder
uses the parameter derivative to detect and refine periodic features
(\cref{sec:feature-finder}); its LA evaluator propagates
\code{dzdc} before updating the delta. The LA helper interfaces also
provide propagation of an incoming-state derivative \code{dzdz}.
The escape-time rendering traversal described above advances deltas
and counts without invoking these derivative helpers.

Given a single LA step with coefficients
\((\code{ZCoeff}, \code{CCoeff})\) and reference \(z_\mathrm{ref}\), the
derivative updates are:
\begin{align}
\code{dzdz}' &= \code{dzdz}
  \;\cdot\; 2\,(\Delta z + z_\mathrm{ref})
  \;\cdot\; \code{ZCoeff},
\label{eq:la-dzdz}
\\
\code{dzdc}' &= \code{dzdc}
  \;\cdot\; 2\,(\Delta z + z_\mathrm{ref})
  \;\cdot\; \code{ZCoeff}
  \;+\; \code{CCoeff} \cdot S,
\label{eq:la-dzdc}
\end{align}
where $S$ is the stored-derivative scale:
$\code{dzdc}_{\mathrm{stored}}=S\,\partial z/\partial c$.
The Feature Finder LA evaluator sets $S=1$ and retains that scale
during its LA pass; the helper accepts an explicit scale to preserve
the caller's representation.
The factor \(2\,(\Delta z + z_\mathrm{ref})\) is the derivative of the
Prepare step, and \code{ZCoeff} propagates through the linearized
iterations.  The updates \eqref{eq:la-dzdz}--\eqref{eq:la-dzdc} are applied \emph{before}
\code{Evaluate} updates \(\Delta z\) itself, so \(\Delta z\) in the
formulas above is the pre-step value.

Writing $\code{dzdz} = g_x + i\,g_y$, $\code{dzdc} = h_x + i\,h_y$,
$\Delta z + z_\mathrm{ref} = q_x + i\,q_y$, and the scaling factor
$S \in \mathbb{R}$, the real-arithmetic updates are:
\begin{align}
g'_x &= (g_x\,q'_x - g_y\,q'_y)\,Z_x
      - (g_x\,q'_y + g_y\,q'_x)\,Z_y, \\
g'_y &= (g_x\,q'_x - g_y\,q'_y)\,Z_y
      + (g_x\,q'_y + g_y\,q'_x)\,Z_x, \\
h'_x &= (h_x\,q'_x - h_y\,q'_y)\,Z_x
      - (h_x\,q'_y + h_y\,q'_x)\,Z_y
      + C_x\,S, \\
h'_y &= (h_x\,q'_x - h_y\,q'_y)\,Z_y
      + (h_x\,q'_y + h_y\,q'_x)\,Z_x
      + C_y\,S,
\end{align}
where $q'_x = 2\,q_x$ and $q'_y = 2\,q_y$ absorb the factor of~$2$.
These expressions expand the same chain-rule complex products into
real and imaginary components.

\subsubsection{Comparison with BLA}
\label{sec:la-vs-bla}

BLA (\cref{sec:bilinear-approx}) and LA~v2 both precompute coefficient
maps, but differ in their input model and hierarchy:
\begin{itemize}
  \item \textbf{Hierarchy structure.}  BLA builds a generic binary tree
        of composed steps at power-of-two granularity (levels
        \(1, 2, 4, 8, \ldots\)), agnostic of the orbit's dynamics.
        LA~v2 segments the orbit at detected return boundaries or
        synthetic target spans, then composes consecutive entries
        into higher stages. Its spans need not be powers of two or
        fixed multiples of a fundamental period.

  \item \textbf{Period awareness.}  BLA has no notion of periodicity;
        it treats the reference orbit as an opaque sequence.  LA~v2
        uses magnitude and threshold-drop heuristics to identify
        useful boundaries near returns to the critical point. Those
        boundaries can exploit periodic structure, while synthetic
        segmentation also permits construction without a detected
        return. Pixel validity is still checked independently.

  \item \textbf{Prepare step.}  BLA applies a purely linear map at every
        step.  LA~v2 applies a nonlinear Prepare step
        (\cref{eq:la-prepare}) that retains the quadratic term at each
        entry's reference point, then linearizes the subsequent
        iterations. Thus its map is linear in the \emph{prepared}
        delta, whereas BLA is linear in the incoming delta itself.
        Both maps still omit nonlinear contributions over their spans.

  \item \textbf{Kernel structure.}  The BLA kernel interleaves single
        perturbation steps with greedy BLA lookups: do one perturbation
        step, scan the hierarchy for the longest valid BLA skip, apply
        it, and repeat. LA~v2 Full rendering first attempts AT, then
        descends through LA stages, and finally performs perturbation
        for the remaining work. A failed block transfers its starting
        position into the finer stage instead of restarting traversal.
\end{itemize}
Long return spans and a usable AT can make LA~v2 effective near a
mini-Mandelbrot copy. The benefit depends on accepted block lengths,
table-construction cost, memory traffic, and the amount of perturbation
work remaining for each pixel. Hierarchy design alone does not
guarantee a particular speedup or accuracy ordering between the methods.

Both perturbation and linear approximation presuppose the existence of
an accurate reference orbit.  The next chapter addresses how
\FractalShark{} computes this orbit at the extreme precisions that
deep-zoom rendering demands.
